\documentclass[11pt]{article}

\usepackage[margin=1in]{geometry}
\usepackage{amsmath,amssymb,amsthm,mathtools}
\usepackage{graphicx}
\usepackage{microtype}
\usepackage{xcolor}
\usepackage[hidelinks]{hyperref}

\newtheorem{theorem}{Theorem}
\newtheorem{corollary}[theorem]{Corollary}
\newtheorem{lemma}[theorem]{Lemma}
\newcommand{\PP}{\mathbb P}
\newcommand{\CC}{\mathbb C}
\newcommand{\RR}{\mathbb R}
\newcommand{\cI}{\mathcal I}
\newcommand{\cP}{\mathcal P}

\title{The rank-one phase-retrieval incidence variety\\
has exactly $2^N$ irreducible components}
\author{Candidate counterexample packet for arXiv:1506.00674}
\date{17 August 2026}

\begin{document}
\maketitle

\begin{center}
\fbox{\parbox{0.91\textwidth}{\textbf{Status.}
Candidate counterexample, likely valid, pending human review.  The source's
incidence variety is not irreducible in general: when all projection ranks are
one it has exactly $2^N$ irreducible components.}}
\end{center}

\section{Source question and scope}

Edidin studies phase retrieval by $N$ orthogonal projections on an
$M$-dimensional real vector space and introduces a complex incidence variety
$\cI_{k_1,\ldots,k_N,M}$ in the proof of generic phase retrieval
\cite{Edidin}.  Proposition 4.3 controls the union of the irreducible
components containing the real locus.  Remark 4.4 then states that the full
incidence variety is not asserted to have the computed dimension because its
irreducibility is not known.  The source passage is on PDF page 7:

\begin{figure}[ht]
  \centering
  \includegraphics[width=0.98\textwidth]{figures/open_problem_crop.png}
  \caption{Proposition 4.3 and Remark 4.4 of arXiv:1506.00674, PDF page 7.}
\end{figure}

The projective factors displayed as $\PP^{N-1}$ in the source formula are a
typographical slip: the surrounding text has $x,y\in\CC^M\setminus\{0\}$,
so the factors used below are $\PP^{M-1}$.

This packet gives a negative answer to the irreducibility uncertainty and an
exact component classification for the all-rank-one case.  It does not alter
the generic phase-retrieval theorem.

\section{Definitions}

Put $V=\CC^M$, $M\geq2$, and let
$b(v,w)=v^{\mathsf T}w$ be the standard nondegenerate symmetric bilinear form.
For $1\leq k\leq M-1$, let
\[
  \cP_k(M)=\{P\in\CC^{M\times M}:P^{\mathsf T}=P,
  \ P^2=P,\ \operatorname{tr}P=k\}.
\]
For ranks $k_1,\ldots,k_N$, define the incidence variety
\[
\cI_{k_1,\ldots,k_N,M}
=\left\{(P_1,\ldots,P_N,[x],[y]):
 b(y,P_i x)=0\quad(1\leq i\leq N)\right\}
\]
inside
$\prod_i\cP_{k_i}(M)\times\PP(V)\times\PP(V)$.

Let
\[
  U=\{[u]\in\PP(V):b(u,u)\ne0\}.
\]
Every $[u]\in U$ determines
\[
  P_u=\frac{uu^{\mathsf T}}{u^{\mathsf T}u}\in\cP_1(M).
\]

\section{Proof intuition}

For a rank-one projection, the incidence equation splits:
\[
 b(y,P_u x)=\frac{b(u,y)b(u,x)}{b(u,u)}=0.
\]
Thus each measured line must be orthogonal either to $x$ or to $y$.
There are $2^N$ ways to make those choices.  Each choice gives a closed
irreducible piece: after fixing the choice, the relevant lines form an open
part of a product of projective bundles over $[x]$ and $[y]$.  Finally,
$x=e_1$, $y=e_2$, and image lines alternating between $\CC e_2$ and
$\CC e_1$ give a point belonging to exactly one chosen piece.  Hence the
pieces are distinct maximal irreducible subsets, not merely a redundant
closed cover.

\section{Exact component theorem}

\begin{lemma}[Rank-one parametrization]\label{lem:param}
The map
\[
 U\longrightarrow\cP_1(M),\qquad [u]\longmapsto P_u,
\]
is an isomorphism of varieties.  Under this isomorphism,
\[
 b(y,P_u x)=0
 \quad\Longleftrightarrow\quad
 b(u,x)=0\ \text{ or }\ b(u,y)=0.
\]
\end{lemma}

\begin{proof}
The displayed formula for $P_u$ is unchanged when $u$ is rescaled, is
symmetric and idempotent, and has image $\CC u$.  Conversely, if
$P\in\cP_1(M)$, symmetry gives
$\ker P=(\operatorname{im}P)^\perp$.  Idempotence gives
$V=\operatorname{im}P\oplus\ker P$, so the image line is nondegenerate.
Writing it as $\CC u$ therefore gives $b(u,u)\ne0$ and forces
$P=P_u$.  The construction and its inverse (taking the image line) are
regular on these varieties.  Direct multiplication gives
\[
 b(y,P_u x)=\frac{b(u,y)b(u,x)}{b(u,u)},
\]
which proves the final equivalence.
\end{proof}

For $a\geq0$, define
\[
 Z_a=\left\{([x],[u_1],\ldots,[u_a])\in
 \PP(V)\times U^a:b(x,u_j)=0\ (1\leq j\leq a)\right\},
\]
with $Z_0=\PP(V)$.

\begin{lemma}\label{lem:Za}
Every $Z_a$ is irreducible and
\[
 \dim Z_a=M-1+a(M-2).
\]
\end{lemma}

\begin{proof}
Over $\PP(V)$, the hyperplanes $x^\perp$ form the fibers of a rank
$M-1$ vector bundle $K$: it is the kernel of the everywhere-surjective
bundle map
\[
 V\otimes\mathcal O_{\PP(V)}\longrightarrow\mathcal O_{\PP(V)}(1),
 \qquad u\longmapsto b(x,u).
\]
The $a$-fold fiber product of $\PP(K)$ over $\PP(V)$ is locally a product
of the irreducible base with $(\PP^{M-2})^a$; hence it is irreducible and
has dimension $M-1+a(M-2)$.  The condition that every $u_j$ be
nonisotropic removes a closed subset.  The resulting open subset is $Z_a$,
and it is nonempty (take $x=e_1$ and all $u_j=e_2$).  Therefore $Z_a$ is
irreducible of the stated dimension.
\end{proof}

\begin{theorem}[All-rank-one component classification]\label{thm:main}
For $M\geq2$ and $N\geq1$, the variety
$\cI_{1,\ldots,1,M}$ has exactly $2^N$ irreducible components.  Every
component has dimension
\[
 2M-2+N(M-2)
 =N(M-1)+2M-2-N.
\]
Moreover, every component contains real points.
\end{theorem}

\begin{proof}
Use Lemma~\ref{lem:param} to replace each $P_i$ by a nonisotropic image
line $[u_i]$.  For a word
$\varepsilon=(\varepsilon_1,\ldots,\varepsilon_N)\in\{X,Y\}^N$, let
$C_\varepsilon$ be the closed subvariety defined by
\[
 \begin{cases}
 b(u_i,x)=0,&\varepsilon_i=X,\\
 b(u_i,y)=0,&\varepsilon_i=Y.
 \end{cases}
\]
The conditions are equivalently $P_i x=0$ and $P_i y=0$, respectively,
so they are closed in the original projection coordinates.  Lemma
\ref{lem:param} gives the finite closed cover
\[
 \cI_{1,\ldots,1,M}=\bigcup_{\varepsilon\in\{X,Y\}^N}C_\varepsilon.
\]

If $a$ entries of $\varepsilon$ equal $X$, then, after reordering the
$u_i$ coordinates,
\[
 C_\varepsilon\cong Z_a\times Z_{N-a}.
\]
Lemma~\ref{lem:Za} shows that $C_\varepsilon$ is irreducible and has
dimension
\[
 [M-1+a(M-2)]+[M-1+(N-a)(M-2)]
 =2M-2+N(M-2).
\]

It remains to show that these irreducible pieces are exactly the components.
Take real vectors $x=e_1$ and $y=e_2$.  For each $i$, choose
\[
 u_i=\begin{cases}
 e_2,&\varepsilon_i=X,\\
 e_1,&\varepsilon_i=Y.
 \end{cases}
\]
Then $b(u_i,u_i)=1$.  If $\varepsilon_i=X$, the line $\CC u_i$ is
orthogonal to $x$ but not to $y$; if $\varepsilon_i=Y$, it is orthogonal
to $y$ but not to $x$.  The resulting real point belongs to
$C_\varepsilon$ and to no $C_\delta$ with $\delta\ne\varepsilon$.
Thus no $C_\varepsilon$ is contained in another one.  A finite closed
cover by irreducible subsets with this noncontainment property is precisely
the irreducible-component decomposition.  The witnesses also prove the last
assertion.
\end{proof}

The dimension in Theorem~\ref{thm:main} is exactly the formula in source
Proposition 4.3,
$\sum_i k_i(M-k_i)+2M-2-N$, specialized to $k_i=1$.  Thus all $2^N$
components containing real points have the same expected dimension.

\begin{corollary}[Any rank-one factor forces reducibility]\label{cor:mixed}
Suppose exactly $r\geq1$ of $k_1,\ldots,k_N$ equal one.  Then
$\cI_{k_1,\ldots,k_N,M}$ has at least $2^r$ irreducible components.  In
particular, it is reducible.
\end{corollary}

\begin{proof}
Applying Lemma~\ref{lem:param} to the $r$ rank-one indices gives a cover by
$2^r$ closed pattern subsets.  Each pattern has a point belonging to no
other pattern: take $x=e_1$, $y=e_2$, choose the rank-one image lines as in
the proof of Theorem~\ref{thm:main}, and for every remaining index choose a
coordinate orthogonal projection of rank $k_i$.  Such a projection is
diagonal, so $e_2^{\mathsf T}P_i e_1=0$.

Choose an irreducible component through each pattern-exclusive point.  An
irreducible subset of a finite union of closed sets lies in one of them.
Hence components chosen from two different exclusive points cannot coincide,
giving at least $2^r$ components.
\end{proof}

\section{Verification and limitations}

The factorization and pattern witnesses were checked independently using
exact rational arithmetic in dimensions $2$ through $6$ and for all patterns
through $N=4$; the reusable checker is included in the packet.  This is only
a sanity check.  The proof of irreducibility of each $C_\varepsilon$ is the
projective-bundle argument in Lemma~\ref{lem:Za}.

The result answers the source's general irreducibility uncertainty negatively
and gives an exact classification when all ranks are one.  It does not
classify the case $2\leq k_i\leq M-1$ for every $i$.  It also does not expose
a gap in the source: Remark 4.4 carefully avoids using irreducibility, and
Proposition 4.3 remains compatible with the component calculation above.

\section{Novelty check and review recommendation}

Before promotion, the run's registry, solution, attempt, and proof-gap indexes
were searched for arXiv:1506.00674 and close phase-retrieval incidence terms.
No exact prior result was present.  Bounded web/arXiv searches used the exact
remark phrase, the title and ``incidence variety irreducible/reducible,'' and
close variants involving orthogonal projections and components.  They found
the source and later phase-retrieval literature but no paper giving this
$2^N$ component classification.  This is evidence of apparent novelty, not
an exhaustive citation review.

Recommendation: send to a human algebraic geometer or phase-retrieval expert.
The principal review points are the variety isomorphism in
Lemma~\ref{lem:param} and the mixed-rank component-count argument in
Corollary~\ref{cor:mixed}.

\begin{thebibliography}{9}
\bibitem{Edidin}
D. Edidin,
\emph{Projections and phase retrieval},
arXiv:1506.00674 (2015), especially Proposition 4.3 and Remark 4.4.

\bibitem{Hartshorne}
R. Hartshorne,
\emph{Algebraic Geometry},
Graduate Texts in Mathematics 52, Springer, 1977, Chapter I.
\end{thebibliography}

\end{document}
